% Standalone plain TeX file using TeXsis
% Revised subsection: "Computing equilibria" (money-financed deficits)
% For insertion into black8_ver7_tom.tex
%
% Cross-references (\Ep{...}, \use{...}, \Fg{...}) will resolve
% correctly once inserted into the parent document.

\input texsis

\subsection{Computing equilibria}

To compute an equilibrium with a constant stream of government
expenditures $g$, write equation \Ep{Gov_bc_Mdef;b} as
$$
  {M_t \over p_t} = {M_{t-1} \over p_{t-1}}\,{p_{t-1} \over p_t}
     + g, \quad \forall\, t \geq 1,
$$
where ${M_0 / p_0}$ is given by equation \Ep{Gov_bc_Mdef;a}.
By substituting the equilibrium condition $M_t / p_t = f(R_t)$, where
$R_t = p_t / p_{t+1}$, into these equations, we arrive at the same
equilibrium characterization as that of an economy with bond-financed
deficits in equations \Ep{DEq2g}.
Thus, the stationary-equilibrium characterization \Ep{DEq2gSS} remains
the same:
$$
  (1-R)\,f(R) = g \quad \Longrightarrow \quad
  {(p_{t+1}/p_t) - 1 \over p_{t+1}/p_t}\,
        f\!\left({p_t \over p_{t+1}}\right) = g,
  \EQN DEq2mSS
$$
where we have used $R_t = p_t / p_{t+1}$ to illuminate the taxation
analogy introduced in subsection \use{sec:SSLaffer}.
Recall that $(1-R)$ was interpreted as a `tax rate'---here
approximately equal to the net inflation rate $(p_{t+1}/p_t) - 1$
at low inflation---and $f(R)$ as a `tax base,' namely agents'
demand for real money balances, which equals the real money supply
$M_t/p_t$ in equilibrium.
Printing money causes inflation and thereby erodes the real value of
money holdings; in a stationary equilibrium agents replenish nominal
money holdings so that real balances remain unchanged.
The government can thus be thought of as levying a `tax' at a rate
approximately equal to the net inflation rate on agents' real money
holdings, the `tax base.'
The `tax rate' cannot exceed $100\,\%$ (a money holder cannot fare
worse than money becoming worthless), so the exact expression in
equation \Ep{DEq2mSS} is
$(1-R) = [(p_{t+1}/p_t) - 1]\,/\,(p_{t+1}/p_t) < 1$.

Next, we revisit the example of logarithmic preferences from
subsection \use{sec:Meqlog}.
Recall difference equation \Ep{samdiff1}, which we now reformulate
using a lag operator,
$$
  \left(1 - {w_2 \over w_1}\,L^{-1}\right) p_t
       = {2 \over w_1}\,M_t,
  \EQN samdiff2
$$
where we now allow for a time-varying forcing term $M_t$.
The solution to this difference equation is
$$
  p_t = {2 \over w_1} \sum_{i=0}^{\infty}
             \left({w_2 \over w_1}\right)^{\!i} M_{t+i}
       + \kappa \left({w_1 \over w_2}\right)^{\!t}.
  \EQN samsol2
$$
If the nominal money supply stays constant, the solution simplifies
to the earlier expression \Ep{samsol1}; otherwise the entire
infinite sequence of nominal money supplies is required.

To illustrate such an equilibrium calculation, we proceed under the
assumption that the economy has converged to a stationary equilibrium.
Given the constant inflation rate $p_{t+1}/p_t = 1/R$ that
characterizes a stationary equilibrium, a natural conjecture is that
the nominal money supply grows at a constant rate,
$M_{t+1} = z M_t$ for some $z > 0$.
According to the first proposition in section \use{sec:olg_optimality},
a necessary and sufficient condition for the existence of a monetary
equilibrium is then $z\,u'(y^i_i)/u'(y^i_{i+1}) < 1$, which here
translates to $z\,w_2/w_1 < 1$.%
\NFootnote{As detailed in section \use{sec:olg_optimality},
Wallace (1980) shows sufficiency of the stated condition by proving
the existence of a stationary monetary equilibrium, which is a
counterpart to the stationary solution of our difference equation
\Ep{samsol2z}.
While Wallace replaces our unfunded government expenditures $g$ with
lump-sum transfers to agents, that difference in assumptions does not
affect Wallace's proofs of necessity and sufficiency for the existence
of at least one monetary equilibrium.
(The difference in assumptions does, of course, matter for any welfare
analysis.)}
Given the sequence
$\{M_{t+i}\}_{i=0}^{\infty} = \{z^i M_t\}_{i=0}^{\infty}$
with $z < w_1/w_2$, we solve equation \Ep{samsol2} to obtain price
sequences consistent with money-market clearing,
$$
  p_t = {2 \over w_1}\,{M_t \over 1 - {w_2 \over w_1}\,z}
       + \kappa \left({w_1 \over w_2}\right)^{\!t}.
  \EQN samsol2z
$$
Setting $\kappa = 0$ yields the price sequence that clears the money
market in a stationary equilibrium in which the nominal money supply
grows at a constant rate $z > 0$.
The associated inflation rate is
$$
  {p_{t+1} \over p_t}
     = { {2 \over w_1}\,{M_{t+1} \over 1 - {w_2 \over w_1}\,z}
         \over
         {2 \over w_1}\,{M_t    \over 1 - {w_2 \over w_1}\,z} }
     = {M_{t+1} \over M_t} = z,
  \EQN DEq2mSSPz
$$
and the real money supply is
$$
  {M_t \over p_t}
     = {M_t \over {2 \over w_1}\,{M_t \over 1 - {w_2 \over w_1}\,z}}
     = {w_1 - w_2 z \over 2} \quad = f(1/z),
  \EQN DEq2mSz
$$
where the last equality is the savings function evaluated at the rate
of return $R = p_t/p_{t+1} = 1/z$.

In a stationary equilibrium with money-financed government deficits
equal to $g$, the growth rate of the nominal money supply $z$ must
satisfy equilibrium condition \Ep{DEq2mSS}:
$$
  {(p_{t+1}/p_t) - 1 \over p_{t+1}/p_t}\,
        f\!\left({p_t \over p_{t+1}}\right) = g
  \quad \Longrightarrow \quad
  {z - 1 \over z}\,f\!\left({1 \over z}\right) = g.
  \EQN DEq2mSSz
$$
As depicted in Figure \Fg{graph52f}, there are two values of $z$
satisfying equation \Ep{DEq2mSSz}: a high-inflation-rate equilibrium
with $z = 1/\underline{R} \gg 1$ and a low-inflation-rate equilibrium
with $z = 1/\overline{R} > 1$.
Drawing on the isomorphism established in subsection
\use{sec:iso_deffin}---that equilibria are identical whether the
government finances its deficits by printing money or by borrowing at
a negative real net interest rate---we can apply the analysis of
bond-financed equilibria in subsection \use{sec:DefBnonstat}.
The low-inflation-rate stationary equilibrium is attained in period $0$
by setting the initial transfer $-\tau^{-1}_0$ equal to the critical
value in equation \Ep{tau_highR}.
If the initial transfer is smaller, however, all monetary equilibria
converge to the high-inflation-rate stationary equilibrium.
As noted earlier, the ``bad'' Laffer-curve equilibrium is stable.

The undesirable stability of the bad Laffer-curve equilibrium is a
feature of perfect foresight dynamics and raises the question of what
alternative solution concepts predict.
Bruno and Fischer (1990) and Marcet and Sargent (1989) analyze how the
system behaves under two different types of adaptive dynamics.
They find that under either a crude form of adaptive expectations or a
least-squares learning scheme, $R_t$ converges to $\overline{R}$.
This finding is reassuring because the comparative dynamics are more
plausible at $\overline{R}$ (larger deficits bring higher inflation).
Furthermore, Marimon and Sunder (1993) present experimental evidence
supporting the equilibrium selection implied by adaptive dynamics.
Marcet and Nicolini (2003) build and calibrate an adaptive model of
several Latin American hyperinflations based on this equilibrium
selection.
Sargent, Williams, and Zha (2009) extend and estimate that model.

\auth{Williams, Noah}%
\auth{Zha, Tao}%
\auth{Marcet, Albert}%
\auth{Nicolini, Juan Pablo}%
\auth{Marimon, Ramon}%
\auth{Sunder, Shyam}%
\auth{Bruno, Michael}%
\auth{Fischer, Stanley}%
\auth{Sargent, Thomas J.}%

\bye
